\section{吞噬边界：Interface、Super App 与 Chatbot 到 Agent}

上一章讲 Agent 内部机制，本章讲外部边界。姚顺雨认为创业公司的最大机会，是设计不同 interface。模型能力可能产生 beyond ChatGPT 的交互方式，变成新的 Super App。这里的“吞噬边界”不是说一个公司会吃掉所有东西，而是说不同交互界面会决定智能系统能进入哪些任务。

\begin{figure}[H]
\centering
\includegraphics[width=0.96\textwidth]{interface-superapp.png}
\caption{吞噬边界：Interface 决定机会，不同交互方式可能产生不同 Super App。自制概念图，依据 00:48:38--01:05:01 对谈内容整理。}
\end{figure}

\begin{knowledgebox}{读图：interface 是边界，不只是 UI}
Chatbot、Assistant、Her、非人形界面、浏览器、IDE、操作系统都可能成为智能入口。不同 interface 会引导不同任务、数据和商业模式，也会塑造研究方向。
\end{knowledgebox}

\subsection{Chatbot 系统为何自然演化成 Agent}

模型公司的 Chatbot 系统天然有用户意图、上下文、工具调用需求和长期状态需求。因此它会自然演化成 Agent 系统：先记住用户，再调用工具，再执行任务，再把反馈写回系统。Chatbot 不是终点，而是 Agent 的入口形态之一。

\begin{figure}[H]
\centering
\includegraphics[width=0.96\textwidth]{chatbot-to-agent.png}
\caption{Chatbot 到 Agent：模型公司的聊天系统会自然演化成 Agent 系统。自制概念图，依据 01:49:28--02:05:00 对谈内容整理。}
\end{figure}

\begin{importantbox}{Chatbot 到 Agent 的最小迁移}
一个 Chatbot 只要加入长期 memory、工具调用、任务状态和结果反馈，就会开始变成 Agent。区别不在界面是否还是聊天框，而在系统是否能持续改变外部状态。
\end{importantbox}

\subsection{Super App 是双刃剑}

上一节讲 Chatbot 如何自然演化成 Agent，本节看强入口带来的组织路径依赖。拥有 ChatGPT 这样的 Super App 是巨大优势，因为它带来用户、数据、分发和产品反馈；但它也会塑造研究方向，让组织自然围绕这个 Super App 优化。创业公司如果只是复制旧 interface，很难超越；如果找到不同 interface，就可能打开新任务边界。

\begin{warningbox}{界面锁定会影响研究路线}
当公司拥有强入口，它的研究很容易服务这个入口。这不是坏事，但会让它低估其他 interface 的机会。历史上的平台替代，常常不是旧平台的更好版本，而是完全不同的交互方式。
\end{warningbox}

\begin{knowledgebox}{Interface 类型表}
\begin{tabularx}{\textwidth}{p{0.22\textwidth}p{0.32\textwidth}X}
\toprule
Interface & 典型任务 & 潜在机会 \\
\midrule
Chatbot & 问答、写作、解释、轻量任务 & 通用入口强，但容易同质化。 \\
IDE/Code & 编程、调试、自动化工具 & affordance 强，任务可验证。 \\
Browser/OS & 跨网页、文件和应用的数字行动 & 可能形成 universal digital agent。 \\
Non-human UI & 非人形、嵌入式、背景式交互 & 可能避开 Chatbot 路径锁定。 \\
\bottomrule
\end{tabularx}
\end{knowledgebox}

因此，“吞噬的边界”可以理解为 interface 的边界。一个模型如果只能在聊天框里工作，它吞噬的是问答、写作和轻量任务；如果它进入 IDE，它吞噬的是代码生产；如果它进入浏览器和操作系统，它吞噬的是跨应用数字行动；如果它进入企业流程，它吞噬的是组织中的重复协调。边界扩张不是模型自己完成的，而是模型、界面、工具和数据共同完成的。

\subsection{本章小结}

智能边界不是由单一模型决定，而是由 interface、任务、工具、用户习惯和组织路径共同决定。Agent 的机会，也会从新的交互方式里长出来。

